\documentclass{article}
%  ╔══════════════════════════════════════════════════════════╗
%  ║                         Title                            ║
%  ╚══════════════════════════════════════════════════════════╝
%NOTE: I want \title,\author to be near top, and this \renewcommand{\maketitle} fails if these are above it.
\usepackage[utf8]{inputenc}         % Used to compile any UTF-8 Character (and supports Unicode)
\usepackage{titling}

\renewcommand{\maketitle}{
    \quad
    \vspace{1em}
  \begin{center}
      \LARGE\bfseries \thetitle \par
      \vspace{0.6em}
      \large
  \end{center}
    \vspace{1em}
}

\title{Geometric Intuition of Hilbert Spaces}
\author{Nathanael Chwojko-Srawley}
\date{\today}

%  ╔══════════════════════════════════════════════════════════╗
%  ║                         packages                         ║
%  ╚══════════════════════════════════════════════════════════╝

\usepackage{extarrows}         % Used to compile any UTF-8 Character (and supports Unicode)
\usepackage{stackengine}
\usepackage{colortbl}
\usepackage{scalerel}
\usepackage{amsmath, amssymb} % Used to write better mathematical expressions (ex. \mathbb{R})
\usepackage{stackrel}         % for left-right arrows, staking symbol above and bellow
\usepackage{mathrsfs}
\usepackage{bbm}                     % for mathbbm{1}
\usepackage{euscript}
\usepackage{commutative-diagrams}
\usepackage[font=itshape]{quoting}
\usepackage{mathtools}
  \usepackage{enumitem}
\usepackage{amsthm}           % Used to create theorem enviroments.
\usepackage{graphicx}         % Let's you use images in LaTeX; ex the \includegraphics command.
\usepackage{hhline}
%\usepackage{luatexja}
%\usepackage{fontspec}
\usepackage{dsfont}
\usepackage{wasysym}          % emoji's!
\usepackage{float}            % package with ``H'' for figures and tables, ex. Images, tables, etc
%\usepackage{subcaption}      % More options with sub-captions.
\usepackage{setspace}         % More flexibility on spacing in document.
\usepackage{hyperref}         % let's me put hyperlinks
%\usepackage{tikz}             % for commutative diagrams https://tex.stackexchange.com/questions/419221/how-to-do-the-pushout-with-universal-property
\usepackage{tikz-cd}          % for commutative diagrams https://tex.stackexchange.com/questions/419221/how-to-do-the-pushout-with-universal-property
%\usepackage{quiver}           % for better commutative diagrams
\usepackage{bookmark}         % creating heading shortcut bar on the left as in word
\usepackage{longtable}        % create tables that extend past one page
%\usepackage{siunitx}         % required for alignment
%\usepackage{multirow}        % for multiple rows
\usepackage{microtype}        % makes nice text. Fixes some under-/over- filled box problems.
%\usepackage{fancyvrb}        % Let's you write code (different style than listings).
\usepackage[most]{tcolorbox} 	      % Makes nice boxes for thms, cor, prop, or anything else.
%\usepackage{enumitem}        % More flexibility on enumeration (don't know much).
\usepackage{xcolor}           % Colour text.
\usepackage{wrapfig}          % To wrap text around figure
\usepackage{listings}         % Create nice colour code blocks (customized in preamble)
\usepackage{thmtools}         % Customize theorem environment (in preamble).
%\usepackage{marvosym}        % Provide ``basic'' symbols (hazard, religious, smiley face etc.)
\usepackage{array}            % \begin{table}{>{\em}c c} -> 1st column italic (bfseries for bold)
%\usepackage[english]{babel}  % If I want to blindtext in english.
\usepackage{blindtext}
\usepackage[a4paper, total={6in, 8.5in}]{geometry}
%\usepackage[symbol]{footmisc} % For daggers and asterisk for footnotes. 1: *, 2: dagger, etc.
\usepackage{fancyhdr}
\usepackage{titlesec}
%\usepackage{pst-plot}          %To make nice plots: https://tex.stackexchange.com/questions/47594/how-can-i-make-a-graph-of-a-function#47596
% \usepackage[answerdelayed]{exercise}
\usepackage{etoolbox}          % for Dynkin diagrams
\usepackage{dynkin-diagrams}   % for Dynkin diagrams
\usepackage{cleveref}

% ╔═════════════════════════════════════════════════════╗
% ║           for figures via inkspace                  ║
% ╚═════════════════════════════════════════════════════╝

% to import figures via: https://github.com/gillescastel/inkscape-figures
\usepackage{import}
\usepackage{pdfpages}
\usepackage{transparent}


% This command is the ONLY command imported via the packages snippet, think of it as a tiny package written out
\newcommand{\incfig}[2][1]{%
    \def\svgwidth{#1\columnwidth}
    \import{./figures/}{#2.pdf_tex}
}

% ╔═══════════════════════════════════════════════════════════╗
% ║                 bibliography and citing                   ║
% ╚═══════════════════════════════════════════════════════════╝

\usepackage{csquotes}
% \usepackage[backend=biber,citestyle=alphabetic,bibstyle=authortitle]{biblatex}
\usepackage{biblatex}

\addbibresource{bibliography.bib}

% ╔═══════════════════════════════════════════════════════════╗
% ║                    colours (colors)                       ║
% ╚═══════════════════════════════════════════════════════════╝

\definecolor{dkgreen}{rgb}{0,0.6,0}
\definecolor{gray}{rgb}{0.5,0.5,0.5}
\definecolor{mauve}{rgb}{0.58,0,0.82}
\definecolor{lightyellow}{rgb}{1,1,0.6}
\definecolor{reddish}{HTML}{A01A3A}

%  ╔══════════════════════════════════════════════════════════╗
%  ║                       book format                        ║
%  ╚══════════════════════════════════════════════════════════╝

\pagestyle{fancy}

\setlength\headheight{20pt}

\renewcommand{\sectionmark}[1]{\markright{\thesection.\ #1}}

\fancyfoot[C,CO]{\textcolor{gray}{Nathanael Chwojko-Srawley}}
\fancyfoot[R,RO]{\thepage}{}

% Create a new page style for the first page
\fancypagestyle{firstpage}{
  \fancyhf{} % Clear header and footer
  \fancyhead[L]{\theauthor}
  \fancyhead[R]{\today} % Date in the top-right
  \fancyfoot[C]{\textcolor{gray}{Nathanael Chwojko-Srawley}}
  \fancyfoot[R]{\thepage}
}

% Apply the first page style conditionally
\AtBeginDocument{\thispagestyle{firstpage}}

%have format the title command
\newcommand{\chapter}[1]{%
  \clearpage
  \vspace*{30pt} % Space before the chapter title
  {\normalfont\huge\bfseries #1\par} % Chapter title formatting
  \vspace{20pt} % Space after the chapter title
}


\setlength{\parindent}{0pt} % don't like indentation infront of paragraphs
\setlength{\parskip}{1ex}   %so that there is still space between paragarphs


%  ╔══════════════════════════════════════════════════════════╗
%  ║                       Shortcuts                          ║
%  ╚══════════════════════════════════════════════════════════╝

%\ExerciseLevelInToc

% I like original mathcal better, but need lowercase.
\DeclareFontFamily{OT1}{pzc}{}
\DeclareFontShape{OT1}{pzc}{m}{it}{<-> s * [1.10] pzcmi7t}{}
\DeclareMathAlphabet{\mathpzc}{OT1}{pzc}{m}{it}


% mathbb
\newcommand{\A}{{\mathbb{A}}}
\newcommand{\C}{{\mathbb{C}}}
\newcommand{\F}{{\mathbb{F}}}
\newcommand{\N}{{\mathbb{N}}}
\newcommand{\Q}{{\mathbb{Q}}}
\newcommand{\R}{{\mathbb{R}}}
\newcommand{\V}{{\mathbb{V}}}
\newcommand{\Z}{{\mathbb{Z}}}
\newcommand{\BI}{{\mathbb{I}}}
\renewcommand{\H}{{\mathbb{H}}}
\renewcommand{\P}{{\mathbb{P}}}


% mathcal
% \newcommand{\co}{{\mathfrak{o}}} %mathcal{o} looks like wreath product, so I changed it to mathfrak
\newcommand{\co}{{\mathpzc{o}}}
\newcommand{\CA}{{\mathcal{A}}}
\newcommand{\CB}{{\mathcal{B}}}
\newcommand{\CC}{{\mathcal{C}}}
\newcommand{\CD}{{\mathcal{D}}}
\newcommand{\CF}{{\mathcal{F}}}
\newcommand{\CG}{{\mathcal{G}}}
\newcommand{\CH}{{\mathcal{H}}}
\newcommand{\CI}{{\mathcal{I}}}
\newcommand{\CK}{{\mathcal{K}}}
\newcommand{\CL}{{\mathcal{L}}}
\newcommand{\CM}{{\mathcal{M}}}
\newcommand{\CN}{{\mathcal{N}}}
\newcommand{\CO}{{\mathcal{O}}}
\newcommand{\CQ}{{\mathcal{Q}}}
\newcommand{\CR}{{\mathcal{R}}}
\newcommand{\CS}{{\mathcal{S}}}
\newcommand{\CT}{{\mathcal{T}}}
\newcommand{\CV}{{\mathcal{V}}}
\newcommand{\CZ}{{\mathcal{Z}}}
% EXCPETION:
\newcommand{\CP}{{\mathcal{P}}}
% \newcommand{\CP}{\mathbb{C}\mathrm{P}}
\newcommand{\RP}{\mathbb{R}\mathrm{P}}


%Euscript
\newcommand{\EB}{\EuScript{B}}
\newcommand{\EC}{{\EuScript{C}}}
\newcommand{\EF}{{\EuScript{F}}}
\newcommand{\EL}{{\EuScript{L}}}
\newcommand{\EO}{{\EuScript{O}}}


%mathfrak (lowercase)
\newcommand{\fa}{{\mathfrak{a}}}
\newcommand{\fb}{{\mathfrak{b}}}
\newcommand{\fc}{{\mathfrak{c}}}
\newcommand{\fd}{{\mathfrak{d}}}
\newcommand{\fm}{{\mathfrak{m}}}
\newcommand{\fn}{{\mathfrak{n}}}
\newcommand{\fo}{{\mathfrak{o}}}
\newcommand{\fp}{{\mathfrak{p}}}
\newcommand{\fq}{{\mathfrak{q}}}


%mathfrak (upper case)
\newcommand{\FA}{{\mathfrak{A}}}
\newcommand{\FB}{{\mathfrak{B}}}
\newcommand{\FC}{{\mathfrak{C}}}
\newcommand{\FD}{{\mathfrak{D}}}
\newcommand{\FK}{{\mathfrak{K}}}
\newcommand{\FM}{{\mathfrak{M}}}
\newcommand{\FN}{{\mathfrak{N}}}
\newcommand{\FO}{{\mathfrak{O}}}
\newcommand{\FP}{{\mathfrak{P}}}


%mathscr
\newcommand{\ff}{{\mathscr{F}}}
\newcommand{\SA}{\mathscr{A}}
\newcommand{\SC}{\mathscr{C}}
\newcommand{\SF}{\mathscr{F}}
\newcommand{\SG}{\mathscr{G}}
\newcommand{\SH}{\mathscr{H}}
\newcommand{\SI}{\mathscr{I}}
\newcommand{\SK}{\mathscr{K}}
\newcommand{\SN}{\mathscr{N}}
\newcommand{\SQ}{\mathscr{Q}}
\newcommand{\SR}{\mathscr{R}}
\newcommand{\SV}{\mathscr{V}}
\newcommand{\SZ}{\mathscr{Z}}


% operatorname -- 2 letters (lowercase and upper case)
\newcommand{\ab}{{\operatorname{ab}}}
\newcommand{\ad}{{\operatorname{ad}}}
\newcommand{\ch}{{\operatorname{ch}}}
\newcommand{\ev}{{\operatorname{ev}}}
\newcommand{\id}{{\operatorname{id}}}
\newcommand{\im}{{\operatorname{im}}}
\newcommand{\nm}{{\operatorname{Nm}}}
\newcommand{\op}{{\operatorname{op}}}
\newcommand{\rk}{{\operatorname{rk}}}
\newcommand{\sh}{{\operatorname{sh}}}
\newcommand{\tr}{{\operatorname{tr}}}

\newcommand{\Ab}{{\operatorname{Ab}}}
\newcommand{\Br}{{\operatorname{Br}}}
\newcommand{\Ch}{{\operatorname{Ch}}}
\newcommand{\Cl}{{\operatorname{Cl}}}
\newcommand{\GL}{{\operatorname{GL}}}
\newcommand{\Nm}{{\operatorname{Nm}}}
\newcommand{\SL}{{\operatorname{SL}}}

\renewcommand{\Re}{{\operatorname{Re}}}
\renewcommand{\Im}{{\operatorname{Im}}}


%categories
\newcommand{\Grp}{{\textbf{Grp}}}
\newcommand{\Fld}{{\textbf{Fld}}}
\newcommand{\Sh}{{\textbf{Sh}}}
\newcommand{\Set}{{\textbf{Set}}}
\newcommand{\Mod}{{\textbf{Mod}}}
\newcommand{\PSh}{{\textbf{PSh}}}
\newcommand{\Sch}{{\textbf{Sch}}}
\newcommand{\Top}{{\textbf{Top}}}
\newcommand{\RgSp}{{\textbf{RgSp}}}
\newcommand{\Ring}{{\textbf{Ring}}}


%operatorname -- 3+ letters (lowercase and upper case)
\newcommand{\rad}{{\operatorname{rad}}}
\newcommand{\sep}{{\operatorname{sep}}}
\newcommand{\res}{{\operatorname{res}}}
\newcommand{\sgn}{{\operatorname{sgn}}}
\newcommand{\pre}{{\operatorname{pre}}}
\newcommand{\ord}{{\operatorname{ord}}}
\newcommand{\vol}{{\operatorname{vol}}}
\newcommand{\lcm}{{\operatorname{lcm}}}

\newcommand{\conj}{{\operatorname{conj}}}
\newcommand{\disc}{{\operatorname{disc}}}
\newcommand{\stab}{{\operatorname{stab}}}
\newcommand{\kdim}{{\operatorname{kdim}}}
\newcommand{\diag}{{\operatorname{diag}}}

\newcommand{\trdeg}{\operatorname{trdeg}}
\newcommand{\coker}{{\operatorname{coker}}}
\newcommand{\codim}{{\operatorname{codim}}}

\newcommand{\Ann}{{\operatorname{Ann}}}
\newcommand{\Aut}{{\operatorname{Aut}}}
\newcommand{\Der}{{\operatorname{Der}}}
\newcommand{\Div}{{\operatorname{Div}}}
\newcommand{\Emb}{{\operatorname{Emb}}}
\newcommand{\End}{{\operatorname{End}}}
\newcommand{\Ext}{{\operatorname{Ext}}}
\newcommand{\Gal}{{\operatorname{Gal}}}
\newcommand{\Hom}{{\operatorname{Hom}}}
\newcommand{\Ind}{{\operatorname{Ind}}}
\newcommand{\Inn}{{\operatorname{Inn}}}
\newcommand{\Int}{{\operatorname{int}}}
\newcommand{\Irr}{{\operatorname{Irr}}}
\newcommand{\Jac}{{\operatorname{Jac}}}
\newcommand{\Mor}[1]{\operatorname{Mor}(\textbf{#1})}
\newcommand{\Nil}{{\operatorname{Nil}}}
\newcommand{\Obj}{{\operatorname{Obj}}}
\newcommand{\Out}{{\operatorname{Out}}}
\newcommand{\Pic}{{\operatorname{Pic}\,}}
\newcommand{\Rep}{{\operatorname{Rep}}}
\newcommand{\Res}{{\operatorname{Res}}}
\newcommand{\Spl}{{\operatorname{Spl}}}
\newcommand{\Syl}{{\operatorname{Syl}}}
\newcommand{\Sym}{{\operatorname{Sym}}}
\newcommand{\Tor}{{\operatorname{Tor}}}

\newcommand{\Char}{{\operatorname{char}}}
\newcommand{\Frac}{{\operatorname{Frac}}}
\newcommand{\Frob}{{\operatorname{Frob}}}
\newcommand{\Proj}{{\operatorname{Proj}\,}}
\newcommand{\RMod}{{}_R\operatorname{Mod}}
\newcommand{\SExt}{{\mathcal{E}\emph{xt}}}
\newcommand{\SHom}{{\emph{Hom}}}
\newcommand{\Span}{{\operatorname{span}}}
\newcommand{\Spec}{{\operatorname{Spec}\,}}
\newcommand{\Supp}{{\operatorname{Supp}}}
\newcommand{\Tens}[1]{#1\otimes --}
\newcommand{\fHom}[1]{\operatorname{Hom}(#1, --)}

\newcommand{\mSpec}{{\operatorname{mSpec}}}
\newcommand{\cofHom}[1]{\operatorname{Hom}(--, #1)}


% connectors
% \renewcommand{\mid}{\big|}
\newcommand{\sse}{\subseteq}
\newcommand{\subg}{\leqslant}
\newcommand{\supg}{\geqslant}
\newcommand{\ssne}{\subsetneq}
\newcommand{\isosubg}{\lesssim}
\newcommand{\nsse}{\not\subseteq}
\newcommand{\nsubg}{\trianglelefteq}
\newcommand{\nsupg}{\trianglerighteq}
\newcommand{\csubg}{\blacktriangleleft}
\renewcommand{\mid}{\big|}


% Arrows
\newcommand{\rw}{\rightarrow}
\newcommand{\Rw}{\Rightarrow}
\newcommand{\lw}{\leftarrow}
\newcommand{\Lw}{\Leftarrow}
\newcommand{\lrw}{\leftrightarrow}
\newcommand{\LRw}{\Leftrightarrow}
\newcommand{\acton}{\curvearrowright}
\newcommand{\actson}{\curvearrowright}
\newcommand\mapsfrom{\mathrel{\reflectbox{\ensuremath{\mapsto}}}}

% arrows for commutative diagram: create four new angled double-struck arrows
\newcommand\myrot[1]{\mathrel{\rotatebox[origin=c]{#1}{$\Rightarrow$}}}
\newcommand\NEarrow{\myrot{45}}
\newcommand\NWarrow{\myrot{135}}
\newcommand\SWarrow{\myrot{-135}}
\newcommand\SEarrow{\myrot{-45}}


%shorthands
\newcommand{\oln}[1]{{\overline{#1}}}
\renewcommand{\bf}[1]{\textbf{#1}}
\newcommand{\qn}{\quad\newline}


% limits
\newcommand{\Lim}{{\varprojlim}}
\newcommand{\coLim}{{\varinjlim}}
\newcommand{\Dlim}{\lim\limits_{\longrightarrow}}
\newcommand{\coDlim}{\lim\limits_{\longleftarrow}}


%for saturated ideals in the localization section (I don't like these, I think I'll delete)
\newcommand{\LIm}{{}^e }
\newcommand{\LPre}{{}^c }

%  ╔══════════════════════════════════════════════════════════╗
%  ║                    Custom Commands                       ║
%  ╚══════════════════════════════════════════════════════════╝

% swap varphi and phi
\let\temp\phi
\let\phi\varphi
\let\varphi\temp

\newcommand{\placeholder}{\_\_\_\_\_}

\newcommand{\set}[2]{\left\{#1 \ \middle|\ #2\right\}}

%mod should not have the crazy amount of space to its right!
\renewcommand{\mod}{\,\operatorname{mod}\,}

% dcups
\newcommand{\dcup}{\sqcup}
\newcommand{\Dcup}{\coprod}
% \newcommand{\Dcup}{\bigsqcup}

\newcommand{\nbot}{\mathrel{\rlap{$\bot$}\mkern2mu/}}

%a nice large circle
\newcommand{\tikzcircle}[2][red,fill=black]{\tikz[baseline=-0.5ex]\draw[#1,radius=#2] (0,0) circle ;}%

\makeatletter
\newcommand*\bigcdot{\mathpalette\bigcdot@{.5}}
\newcommand*\bigcdot@[2]{\mathbin{\vcenter{\hbox{\scalebox{#2}{$\m@th#1\bullet$}}}}}
\makeatother

%somehow not a default command
\def\rddots{\mathstrut^{.^{.^{.}}}}

%% new delimiter
\DeclarePairedDelimiter{\ceil}{\lceil}{\rceil}


%  ╔══════════════════════════════════════════════════════════╗
%  ║                     environments                         ║
%  ╚══════════════════════════════════════════════════════════╝


%remember the option: breakable

% create tcolor boxes
\newtcbtheorem[number within=section]{thm}{Theorem}%
{
colback=green!5,
colframe=green!35!black,
fonttitle=\bfseries,
label type=theorem
}{th}

\newtcbtheorem[number within=section, use counter from=thm]{lem}{Lemma}%
{
colback=blue!5,
colframe=black!35!black,
fonttitle=\bfseries,
label type=lemma
}{lm}
\newtcbtheorem[number within=section, use counter from=thm]{prop}{Proposition}%
{
colback=white!5,
colframe=white!35!black,
fonttitle=\bfseries,
label type=proposition
}{pr}
\newtcbtheorem[number within=section, use counter from=thm]{cor}{Corollary}%
{colback=blue!5,
colframe=blue!35!black,
fonttitle=\bfseries,
label type=corollary
}{co}
\newtcbtheorem[number within=section, use counter from=thm]{axiom}{Axiom}%
{colback=black!5,
colframe=yellow!35!black,
fonttitle=\bfseries,
label type=axiom
}{ax}
\newtcbtheorem[number within=section, use counter from=thm]{defn}{Definition}%
{colback=red!5,
colframe=red!35!black,
fonttitle=\bfseries,
label type=definition
}{df}


% for <s-cr> to create \item[$\LRw$] instead of \item
\newenvironment{equivEnumerate}
 {\begin{enumerate}
	}{
\end{enumerate}}



%Custom enviroments
 \newenvironment{note}
 {\begin{tcolorbox}[enhanced, sharp corners, colback=white , borderline={0.2pt}{0pt}{black}]
   \begin{quote}\textbf{Note}:
   }{
   \end{quote}\end{tcolorbox}}

 \newenvironment{intuition}
 {\begin{quote}\textbf{Intuition}:
   }{
   \end{quote}}

\newtcbtheorem[number within=section, use counter from=thm]{titledBox}{}{%
  enhanced,
  sharp corners,
  colback=white,
  colframe=black,
  borderline={0.2pt}{0pt}{black},
  left=2mm,
  right=2mm,
  top=2mm,
  bottom=2mm,
  title style={opacity=1},
  attach title to upper,
  before upper={\textbf{\tcbtitletext}\quad},
  coltitle=black,
  fonttitle=\bfseries,
  separator sign={\quad}
}{box}


%better example and proof environments
\def\exampletext{Example} % If English
\colorlet{colexam}{red!55!black} % Global example color


\newtcbtheorem[number within=section, use counter from=thm]{example}{Example}{% \newtcolorbox[use counter=testexample]{testexamplebox}{%
% Example Frame Start
empty,% Empty previously set parameters
title={\exampletext \thetcbcounter #1},% use \thetcbcounter to access the testexample counter text
% Attaching a box requires an overlay
attach boxed title to top left,
   % Ensures proper line breaking in longer titles
   minipage boxed title,
% (boxed title style requires an overlay)
boxed title style={empty,size=minimal,toprule=0pt,top=4pt,left=3mm,overlay={}},
coltitle=colexam,fonttitle=\bfseries,
before=\par\medskip\noindent,parbox=false,boxsep=0pt,left=3mm,right=0mm,top=2pt,breakable,pad at break=0mm,
   before upper=\csname @totalleftmargin\endcsname0pt, % Use instead of parbox=true. This ensures parskip is inherited by box.
% Handles box when it exists on one page only
overlay unbroken={\draw[colexam,line width=.5pt] ([xshift=-0pt]title.north west) -- ([xshift=-0pt]frame.south west); },
% Handles multipage box: first page
overlay first={\draw[colexam,line width=.5pt] ([xshift=-0pt]title.north west) -- ([xshift=-0pt]frame.south west); },
% Handles multipage box: middle page
overlay middle={\draw[colexam,line width=.5pt] ([xshift=-0pt]frame.north west) -- ([xshift=-0pt]frame.south west); },
% Handles multipage box: last page
overlay last={\draw[colexam,line width=.5pt] ([xshift=-0pt]frame.north west) -- ([xshift=-0pt]frame.south west); }%
}{ex}



\def\prooftext{Proof} % If English
\NewDocumentEnvironment{Proof}{ O{} }
{
    \colorlet{colexam}{black!55!black} % Global example color
    \newtcolorbox[number within=section]{testproofbox}{%
        empty,% Empty previously set parameters
        title={\emph{\prooftext} \thetcbcounter: #1},
        attach boxed title to top left,
        minipage boxed title,
        boxed title style={empty,size=minimal,toprule=0pt,top=4pt,left=3mm,overlay={}},
        coltitle=colexam,
        fonttitle=\bfseries,
        before=\par\medskip\noindent,
        parbox=false,
        boxsep=0pt,
        left=3mm,
        right=0mm,
        top=2pt,
        breakable,
        pad at break=0mm,
        before upper=\csname @totalleftmargin\endcsname0pt,
        % Removed height constraints
        % Added flexible width
        width=\dimexpr\textwidth-3mm\relax,
        % Modified overlays
        overlay unbroken={\draw[colexam,line width=.5pt] ([xshift=-0pt]title.north west) -- ([xshift=-0pt]frame.south west); },
        overlay first={\draw[colexam,line width=.5pt] ([xshift=-0pt]title.north west) -- ([xshift=-0pt]frame.south west); },
        overlay middle={\draw[colexam,line width=.5pt] ([xshift=-0pt]frame.north west) -- ([xshift=-0pt]frame.south west); },
        overlay last={\draw[colexam,line width=.5pt] ([xshift=-0pt]frame.north west) -- ([xshift=-0pt]frame.south west); },%
    }
    \begin{testproofbox}
}
{\end{testproofbox}\endlist}


%  ╔══════════════════════════════════════════════════════════╗
%  ║                    for Dynkin diagrams                   ║
%  ╚══════════════════════════════════════════════════════════╝

\def\row#1/#2!{#1_{\IfStrEq{#2}{}{n}{#2}} & \dynkin{#1}{#2}\\}
\newcommand{\tble}[1]{
   \renewcommand*\do[1]{\row##1!}
   \[
      \begin{array}{ll}\docsvlist{#1}\end{array}
   \]
}

%  ╔══════════════════════════════════════════════════════════╗
%  ║                         links                            ║
%  ╚══════════════════════════════════════════════════════════╝

\definecolor{LinkColour}{HTML}{A64A3B} % favourite
% \definecolor{LinkColour}{HTML}{8C3D32} % 2nd place
\hypersetup{
  colorlinks,
  citecolor=black,
  filecolor=magenta,
  linkcolor=LinkColour,
  urlcolor=blue,
}


%  ╔══════════════════════════════════════════════════════════╗
%  ║                          misc                            ║
%  ╚══════════════════════════════════════════════════════════╝

\stackMath %to be able to stack text in math mode
\makeindex

%import options: all, bibliography, book-formatting, booksSetting, customEnvironments, dynkin, hw-formatting, language-setup, newcommands, packages, preamble, proofs, settings, tcolorbox, test.svg,
\begin{document}
\maketitle
Banach spaces provide a robust framework for measuring the magnitude of elements, allowing us to quantify
convergence and bound operations via the triangle and Minkowski inequalities. However, a fundamental
deficiency in general Banach spaces is the lack of a linear measurement theory to quantify \emph{similarity}.
While we can measure the total size of a sum $\|f+g\|$, we often lack the granular information required to
isolate the contribution of $f$ from $g$. For instance, determining how much a vector projects onto a specific
basis element is straightforward in Euclidean space but algebraically obstructed in general norms.

To make this deficiency precise, consider how the inner product arises in Euclidean geometry. Define the
\emph{energy} of a vector as $E(x) = \frac{1}{2}\|x\|_2^2$. Then the directional derivative of this energy
at $x$ in the direction of a perturbation $y$ is:
\[
	\frac{d}{dt} E(x + ty) \Big|_{t=0}
		= \frac{d}{dt} \frac{1}{2} \langle x+ty, x+ty \rangle \Big|_{t=0}
		= \langle x, y \rangle
\]
Thus, the inner product $\langle x, y \rangle$ is not just an algebraic artifact; it physically represents
the \emph{sensitivity} of the energy of $x$ to a perturbation by $y$. In particular, the energy satisfies the
clean decomposition:
\[
	E(x + y) = E(x) + E(y) + \langle x, y \rangle
\]
The interaction term $\langle x, y \rangle$ is \emph{bilinear}---it decomposes the interference between two
vectors according to the simple rules of superposition. In general Banach spaces, however, such a clean, linear
decomposition of the interaction is impossible due to the non-flat geometry of the unit ball.

The natural question, then, is: \emph{what replaces the inner product in $L^p$?} The correct generalization is
to replace the quadratic energy with the \emph{$p$-energy}:
\[
	F(f) = \frac{1}{p}\|f\|_p^p = \frac{1}{p}\int |f|^p
\]
and ask: how does $F(f)$ change if we perturb $f$ by a function $g$? To answer this, we must compute the
gradient $\nabla F$, which requires first understanding the natural duality of $L^p$ spaces.

\medskip\noindent\textbf{The Duality Map.}\quad
What does it mean to go in the ``same direction'' in a function space? (For simplicity we work throughout with
$L^p$ spaces for $1 < p < \infty$.\footnote{The reasoning can extend to nuclear Fr\'echet spaces, which admit
descriptions as projective limits of Hilbert spaces. For general metric spaces, a different approach is needed
via CAT($k$) geometry; see~\cite{millerGEOMETRYMEASURECATk} for details.}) Rigorously, we are asking for the
functional $\phi_f$ of unit norm that \emph{maximizes} the evaluation at $f$. We can interpret this $\phi_f$
as the \emph{optimal filter} or sensor for $f$---it captures the full ``energy'' of $f$ without
loss ($\phi_f(f) = \|f\|_p$ and $\|\phi_f\|_p = 1$).

By~\cref{th:descriptionDualLp} (and the equality condition of H\"older's inequality), this functional
corresponds to $g \in L^q$ such that:
\[
	g(x) = |f(x)|^{p-1}\sgn(f(x))
\]
Note that $g \in L^q$ since:
\[
	|g|^q = |f|^{q(p-1)} = |f|^{qp - q} \overset{!}{=}|f|^p
\]
where $\overset{!}{=}$ holds since $1/p + 1/q = 1$ iff $p+q = pq$ iff $p = q(p-1)$. Hence as $f \in L^p$:
\[
	\int |g|^q = \int |f|^p < \infty
\]
Thus, we have a map $f \mapsto \phi_f$. We define the \emph{duality map}:
\[
	J_p: L^p \to L^q, \qquad f \mapsto \phi_f
\]
As $(L^q)^* \cong L^p$, we can equivalently view this as a map $J_p: L^p \to (L^p)^*$.

\medskip\noindent\textbf{The Central Geometric Fact.}\quad
The duality map $J_p$ is not an algebraic accident---it is the \emph{geometric normal} to the level sets of
the $p$-energy. To see this, consider the function:
\[
F(x) = \frac{1}{p} \|x\|_p^p = \frac{1}{p} \sum |x_i|^p
\]
(where the $\ell^p$ norm is used for visualization, but the logic holds for integrals).
From multivariable calculus, the unit sphere is the level set $F(x) = \frac{1}{p}$, and the gradient
$\nabla F(x)$ is the vector strictly perpendicular (normal) to the level set at point $x$.
Computing the gradient (the Fr\'echet derivative):
\[
\frac{\partial}{\partial x_i} \left( \frac{1}{p} |x_i|^p \right) = |x_i|^{p-1} \text{sgn}(x_i)
\]
So, the gradient vector is:
\[
\nabla F(x) = \left( |x_1|^{p-1}\text{sgn}(x_1), \dots, |x_n|^{p-1}\text{sgn}(x_n) \right)
\]
This is exactly our duality map:
\[
	\boxed{\nabla F(x) = J_p(x)}
\]
This identification is the key to the entire theory. It tells us that the optimal functional for
$f$---the element of $(L^p)^*$ that best ``measures'' $f$---is precisely the gradient of the energy at $f$:
the direction of steepest increase in $p$-energy. The duality map is a \emph{geometric} object, not merely an
algebraic one.

\medskip\noindent\textbf{The $p$-Pairing and Its Interpretation.}\quad
With the gradient in hand, we can now answer our original question: how does the $p$-energy $F(f)$ change under
a perturbation by $g$? The first-order variation is\footnote{The $dF$ here can be thought of as the
G\^ateaux derivative, commonly defined for functions between locally convex topological vector spaces.}:
\[
d F(f; g) = \langle \nabla F(f), g \rangle_{\text{dual}} = \int \nabla F(f) \cdot g
\]
Substituting $\nabla F = J_p(f)$, this defines a natural pairing:
\[
	\langle f,g \rangle_p := \int J_p(f) \cdot g
\]
Note that it is important to pass through $J_p(f)$ since $\int fg$ need not be finite in general. However, by
H\"older's inequality:
\[
	\int |J_p(f) \cdot g| < \infty
\]
This pairing has a precise geometric interpretation. The gradient $J_p(f)$ is a cotangent vector---a 1-form
$dF|_f \in T_f^* L^p$---and the perturbation $g$ is a tangent vector $g \in T_f L^p$. Thus the pairing is the
evaluation of a 1-form on a tangent vector:
\[
	\langle f, g \rangle_p = dF|_f(g) \;\in\; T_f^* L^p \times T_f L^p \to \R
\]
This canonical evaluation pairing between the tangent and cotangent spaces exists in \emph{every} geometric
setting without any choice of metric. What $J_p$ adds is the ability to \emph{produce} a specific cotangent
vector from a tangent vector, turning the pairing from a bilinear form on $T_f L^p \times T_f^* L^p$ into a
(generally nonlinear) form on $T_f L^p \times T_f L^p$---that is, internalizing the measurement within the
space itself.

\medskip\noindent\textbf{The Infinitesimal Meaning of $V^*$.}\quad
We are now in a position to make a precise statement about the geometric content of the dual space. In a bare
vector space without additional structure, the dual $V^*$ is merely an algebraic object---a collection of
linear maps with no geometric significance beyond probing linearly. However, in $L^p$ equipped with its norm,
the identity $J_p = \nabla F$ reveals that \emph{every element of $(L^p)^*$ arises as the gradient of the
energy functional}---a differential $dF|_f$. This endows $(L^p)^*$ with genuine infinitesimal meaning: its
elements are not just linear maps, but \emph{cotangent vectors} that encode the first-order sensitivity of the
geometry.

Crucially, cotangent vectors are a different species of infinitesimal object than tangent vectors. A tangent
vector $g \in T_f L^p \cong L^p$ represents an \emph{infinitesimal direction of motion}---it answers ``which
way?''. A cotangent vector $\phi \in T_f^* L^p \cong L^q$ represents an \emph{infinitesimal
measurement}---it answers ``how fast is this quantity changing?''. These are related by the evaluation pairing
$\phi(g) = dF|_f(g)$, but they transform in opposite ways: tangent vectors push forward under maps, while
cotangent vectors pull back. The duality map $J_p: T_f L^p \to T_f^* L^p$ provides a (nonlinear)
identification between these two types of infinitesimal data, allowing us to transfer the directional
interpretation onto the dual space. However, this transfer \emph{depends essentially on the norm
geometry}---it is canonical only for $p = 2$.

Geometrically, if we view the identity map $f \mapsto f$ as the radial vector
field (a \emph{section} of the tangent bundle $TL^p$), then the duality map $J_p$ transforms this into a
covector field (a \emph{section} of the cotangent bundle $T^*L^p$). In the language of Riemannian geometry,
this transformation is the \emph{musical isomorphism} $\flat$ (flat), which maps a vector field $X$ to a
1-form $X^\flat$ by lowering indices via the metric tensor. The four corollaries that follow are all
consequences of this single geometric fact.

\medskip\noindent\textbf{Corollary 1: Nonlinearity of the Musical Isomorphism.}\quad
The duality map is \emph{not} linear:
\[
	J_p(2f) = |2f|^{p-1}\sgn(2f) = 2^{p-1}|f|^{p-1}\sgn(f) = 2^{p-1}J_p(f)
\]
This is only linear in the special case of $p =2$, since:
\[
	J_2(nf) = |nf|^{2-1}\sgn(nf) = n J_2(f)
\]
This distinction is profound. In a general Banach space ($p \neq 2$), the musical isomorphism is
\emph{nonlinear}: the operation of lowering an index depends on the magnitude of the vector itself, meaning
the ``metric'' changes as we move through the space. However, when $p=2$, this map is linear. This implies a
linear bundle isomorphism $TL^2 \cong T^*L^2$, effectively trivializing the distinction between vectors and
covectors. Only in $L^2$ is the musical isomorphism a linear operator, meaning the ``force'' required to
correct an error is linearly proportional to the error itself.

\medskip\noindent\textbf{Corollary 2: Non-Symmetry of the Pairing.}\quad
The pairing $\langle f, g \rangle_p$ is \emph{not} symmetric for $p \neq 2$. This is an immediate consequence
of the nonlinearity of $J_p$: the 1-form $dF|_f$ evaluated on $g$ has no reason to equal the 1-form $dF|_g$
evaluated on $f$, since the gradient depends nonlinearly on its base point. If however $p = 2$, then
$J_p(x)$ maps identically to itself and we recover the symmetric form $\int fg$. This asymmetry is
the first signal that the geometry of $L^p$ is \emph{not flat} for $p \ne 2$.

\medskip\noindent\textbf{Corollary 3: Failure of Rotational Invariance.}\quad
This lack of symmetry in the pairing $\langle f,g \rangle_p$ is directly linked to the fact that $L^2$ is the
unique $L^p$ space that is \emph{rotationally invariant}. Geometrically, rotational invariance means that the
properties of a vector depend only on its length, not on its alignment with the coordinate axes.
Mathematically, this requires that the duality map commutes with rotations. If $R$ is a rotation operator
(an isometry), we expect the ``normal vector'' of the rotated point to be the rotated ``normal
vector'' of the original point:
\[
	J_p(Rf) \overset{?}{=} R(J_p(f))
\]
In $L^2$, since $J_2$ is linear (essentially the identity), this holds trivially: $J_2(Rf) = Rf = R(J_2(f))$.
However, for $p \ne 2$, the map $f \mapsto |f|^{p-1}\sgn(f)$ applies a non-linear distortion component-wise.
Rotating the vector mixes the components, and applying the power law \emph{after} mixing is not the same as
applying it \emph{before}.

To see this concretely, consider $\ell^3$ ($p=3$) in $\R^2$ and let $f = (1, 0)$. If we first compute the functional and then rotate by $\pi/4$ ($45^\circ$), we get:
\[
	R(J_3(f)) = R(1^2, 0^2) = R(1, 0) = \left( \frac{1}{\sqrt{2}}, \frac{1}{\sqrt{2}} \right)
\]
However, if we rotate the vector first, the components become $(\frac{1}{\sqrt{2}}, \frac{1}{\sqrt{2}})$.
Applying the duality map $J_3$ (squaring components) yields:
\[
	J_3(Rf) = J_3\left( \frac{1}{\sqrt{2}}, \frac{1}{\sqrt{2}} \right) = \left( \frac{1}{2}, \frac{1}{2} \right)
\]
Since $(\frac{1}{\sqrt{2}}, \frac{1}{\sqrt{2}}) \ne (\frac{1}{2}, \frac{1}{2})$, we have $J_p(Rf) \ne
R(J_p(f))$. This confirms that for $p \ne 2$, the geometry is \emph{anisotropic}: the shape of the unit ball looks different depending on the angle from which you view it.

\medskip\noindent\textbf{Corollary 4: Birkhoff-James Orthogonality.}\quad
With this understanding of the position-dependent metric, we can clarify the notion of orthogonality.
In Euclidean geometry, orthogonality is symmetric ($x \perp y \iff y \perp x$) and additive. In
general Banach spaces, without a global bilinear form, we must rely on \emph{Birkhoff-James orthogonality}: we say $x \perp y$ if moving in the direction of $y$ does not decrease the length of $x$ ($\|x + \lambda y\| \ge \|x\|$). Geometrically, this means the line passing through $x$ in the direction of $y$ is tangent to the unit sphere at $x$.

Analytically, utilizing the duality map $J_p$, this condition is equivalent to requiring that the tangent
vector $y$ be annihilated by the cotangent vector generated by $x$:
\[
	\langle J_p(x), y \rangle = 0
\]
This is precisely the statement that $y$ lies in the kernel of the 1-form $dF|_x$---a natural geometric
condition. However, it reveals the fundamental asymmetry in non-Hilbert geometry. In $L^2$, because $J_2$ is linear, the condition is symmetric. However, in $L^p$ ($p \ne 2$), the non-linearity of $J_p$ breaks reciprocity. The set of vectors orthogonal to $x$ (kernel of $J_p(x)$) does not align symmetrically with the set of vectors orthogonal to $y$. Consequently, it is common to have $x \perp y$ but $y \not\perp x$. This lack of symmetry illustrates why Hilbert spaces are the unique setting where Euclidean geometric intuition remains valid.

\medskip\noindent\textbf{The Riemannian Perspective.}\quad
The corollaries above all concern the first derivative of the energy (the gradient $J_p$). The full geometry of
the space, however, is determined by how this gradient \emph{changes}---the second derivative,
which acts as the metric tensor. In this framework, the metric tensor $g_{ij}(x)$ at a point $x$ dictates how
to measure the length of tangent vectors. In our context, this metric is generated by the Hessian of the
energy function $F(x) = \frac{1}{p}\|x\|_p^p$. For $L^2$, this Hessian is proportional to the identity
operator everywhere; the metric is \emph{constant} (flat). However, for $L^p$ ($p \ne 2$), the second
derivative depends on position $x$, specifically scaling as $|x|^{p-2}$. This implies that the ``metric''
changes as we move around the space.

It is worth noting why the theory breaks down at the boundary cases. For $L^1$, the unit ball is a diamond
with corners; at these corners, the gradient is not unique (it is a set-valued subgradient). For $L^\infty$,
the unit ball has flat faces; on these faces, the gradient is non-unique or vanishes in certain directions.
This lack of smoothness ($L^1$) or rotundity ($L^\infty$) prevents a well-defined duality map, leaving $L^p$
($1<p<\infty$) as the sweet spot for smooth geometry.

\medskip\noindent\textbf{Topological vs.\ Geometric Structure.}\quad
To visualize the practical consequences of this non-flat geometry, it is helpful to consider the ``shape'' of
the unit ball in infinite dimensions. In finite dimensions, all norms are equivalent, meaning the unit ball of
one norm can always be fit inside a scaled unit ball of another.
In infinite dimensions, this equivalence breaks. The curvature of the norm dictates how ``energy'' is allowed to
concentrate without bound. The $L^\infty$ unit ball acts like a strict box, while the $L^1$ unit ball allows for infinitely long thin spikes.

Ideally, we would like to know if these distortions are topological or ``merely'' geometric. This is answered by the \emph{Mazur Map} $M_{p,q}: L^p(B) \to L^q(B)$ where $B$ is the unit ball. The map:
\[
M_{p,q}(f) = \text{sgn}(f) \cdot |f|^{p/q}
\]
shows that all $L^p(B)$ spaces are \emph{uniformly homeomorphic} to each other (not just homeomorphic
as~\cref{th:andersonKadec} shows). This means that we can smoothly deform the ``spikes'' of $L^1$ into the
``sphere'' of $L^2$. However, while the Mazur map preserves the topology, it destroys the linear structure
(additivity). The Hilbert space $L^2$ is the unique fixed point where this topological structure aligns
perfectly with the linear algebraic structure (where the Mazur map would be the identity).

This geometric disparity explains why convergence depends on the norm. Convergence $f_n \to 0$ is not about
movement, but about entering arbitrarily small $\epsilon$-balls. Consider the sequence $f_n(x) = x^n$ on
$[0,1]$. In the $L^1$ geometry, this sequence travels down one of the ``spikes'' of the unit ball: the functions get
thinner and thinner, their area shrinking to zero despite their height
remaining constant. In the $L^\infty$ geometry, however, these functions are trapped on the surface of the unit
sphere; they cannot enter the $\epsilon$-ball because the $L^\infty$ ``box'' does not extend down
the ``spikes'' of concentration.

\medskip\noindent\textbf{Consequences for Hilbert Spaces.}\quad
The unique geometry of Hilbert spaces---specifically self-duality, isotropy, and the inner product---unlocks
powerful results unavailable in general Banach spaces. The preceding analysis explains \emph{why}: in $L^2$,
the duality map $J_2$ is the identity, making the musical isomorphism linear, the pairing symmetric, and the
geometry flat. The following are key consequences we will explore:

\begin{enumerate}
    \item \emph{The Projection Theorem (Best Approximation):}
    In general Banach spaces, finding the element in a subspace $M$ closest to a point $x$ is difficult (the
    closest point might not be unique, e.g., the flat edge of $L^1$). In Hilbert spaces, the strict convexity and ``roundness'' of the unit ball guarantee that for any closed subspace $M$, there exists a \emph{unique} element $P_M x$ minimizing the distance $\|x - m\|$. Furthermore, the error vector $x - P_M x$ is orthogonal to the subspace.

    \item \emph{The Riesz Representation Theorem (Self-Duality):}
    The duality map $J: H \to H^*$ is linear and bijective. This means we can completely identify linear functionals with vectors. Every continuous linear functional $\phi$ acts exactly like an inner product: $\phi(x) = \langle x, y \rangle$ for a unique $y$. This collapses the study of $H^*$ to the study of $H$.

	\item \emph{Orthonormal Bases and Fourier Series:} In $L^p$, we may have a Schauder basis\footnote{A sequence $\{b_n\}$ where every $v \in V$ has a unique expansion $v = \sum \alpha_nb_n$.}, but determining the coefficients is non-trivial. In Hilbert spaces, orthogonality allows us to decouple coefficients. This is key for the Fourier transform, where any element is reconstructed by projection: $x = \sum \langle x, e_n \rangle e_n$. Additionally, the \emph{Parseval Identity} allows us to measure size by summing squared coefficients: $\|x\|^2 = \sum |\langle x, e_n \rangle|^2$.

    \item \emph{Existence of Adjoints:}
    In Banach spaces, the adjoint of $T: X \to Y$ maps $Y^* \to X^*$. Because Hilbert spaces are self-dual, we can pull this operator back to the original space. Thus, for every bounded operator $T$ on a Hilbert space, there is a unique adjoint $T^*$ on the \emph{same space} satisfying $\langle Tx, y \rangle = \langle x, T^* y \rangle$. This is foundational for spectral theory.
\end{enumerate}

In summary, while Banach spaces provide the necessary topological framework (completeness), Hilbert spaces add
the rigid geometric structure required to generalize linear algebra to infinite dimensions. The duality map
$J_p$, understood as the gradient of the $p$-energy, reveals the precise mechanism: it is the \emph{cotangent
structure} of the space---the infinitesimal measurement theory---that determines the geometry. The canonical
evaluation pairing between tangent and cotangent vectors is the universal source of ``derivative-like''
phenomena; $J_p$ is the geometric data that internalizes this pairing within the space itself. In $L^2$ alone,
this internalization is linear, making vectors and covectors---directions and measurements---two aspects of a
single object.

% \newpage
% \printbibliography

\end{document}
